\begin{proposition}\label{sga3:XXI:6.8.7}
In the notation of \ref{sga3:XXI:6.8.1}, \(u\) induces an
isomorphism from the set of systems of simple roots (respectively
positive roots) of \(R\) onto the corresponding set for \(R'\).
\end{proposition}

\begin{proof}
This follows immediately from \ref{sga3:XXI:3.1.5}
(respectively \ref{sga3:XXI:3.2.1}).
\end{proof}

\section{Structure}\label{sga3:XXI:7}

\subsection{Decomposition of a root datum}\label{sga3:XXI:7.1}

\begin{proposition}\label{sga3:XXI:7.1.1}
Let \(\cR\) be a root datum and let \(\Delta\) be a system of
simple roots.
\begin{enumerate}
\item[\SGAInlineRefTarget{sga3:XXI:7.1.1:item:i}(i)] Let \(R'\) and \(R''\) be two closed symmetric sets of
  roots forming a partition of \(R\). If
  \[
   \Delta'=\Delta\cap R',
   \qquad
   \Delta''=\Delta\cap R'',
  \]
  then
  \[
   R'=R_{\Delta'},
   \qquad
   R''=R_{\Delta''},
  \]
  and every root of \(\Delta'\) is orthogonal to every root of
  \(\Delta''\).
\item[\SGAInlineRefTarget{sga3:XXI:7.1.1:item:ii}(ii)] Let \(\Delta'\) and \(\Delta''\) be two orthogonal
  subsets forming a partition of \(\Delta\). Then
  \(R_{\Delta'}\) and \(R_{\Delta''}\) form a partition of \(R\).
\end{enumerate}
\end{proposition}

\begin{proof}
We first prove \SGARefLink{sga3:XXI:7.1.1:item:i}{(i)}.

\begin{lemma}\label{sga3:XXI:7.1.2}
Under the hypotheses of \textup{\SGARefLink{sga3:XXI:7.1.1:item:i}{(i)}}, if
\(\alpha\), \(\beta\), and \(\alpha+\beta\) are roots, then all
three belong to \(R'\), or all three belong to \(R''\).
\end{lemma}

Suppose, for example, that \(\alpha+\beta\in R'\). One cannot have
\(\alpha,\beta\in R''\), since \(R''\) is closed. Suppose therefore
that \(\alpha\in R'\). Then \(-\alpha\in R'\), and
\[
 \beta=(\beta+\alpha)-\alpha\in R'.
\]

We now prove \(R'=R_{\Delta'}\) by induction on the order of a
positive root \(\alpha\in R'\cap\mathscr P(\Delta)\). If
\(\operatorname{ord}_\Delta(\alpha)=1\), then
\[
 \alpha\in R'\cap\Delta=\Delta'.
\]
If \(\operatorname{ord}_\Delta(\alpha)>1\), there is
\(\beta\in\Delta\) such that \(\alpha-\beta\in R\). By the lemma,
\[
 \beta\in\Delta',
 \qquad
 \alpha-\beta\in R'.
\]
Hence \(\alpha-\beta\in R_{\Delta'}\) by induction, and finally
\[
 \alpha=(\alpha-\beta)+\beta\in R_{\Delta'}.
\]

We next prove that \(\Delta'\) and \(\Delta''\) are orthogonal. If
\(\alpha\in\Delta'\) and \(\beta\in\Delta''\), then
\[
 (\beta^\ast,\alpha)\leq0.
\]
If \((\beta^\ast,\alpha)\neq0\), then \(\beta+\alpha\) is a root,
contrary to the lemma.

We now prove \SGARefLink{sga3:XXI:7.1.1:item:ii}{(ii)}. If \(\Delta'\) or \(\Delta''\) is empty, the
assertion is immediate. Otherwise, if \(R_{\Delta'}\) and
\(R_{\Delta''}\) do not form a partition of \(R\), there is a root
\(\alpha\) of the form
\[
 \alpha=\sum_i m_i'\alpha_i'
       +\sum_j m_j''\alpha_j'',
 \qquad
 m_i',m_j''\in\mathbb Z_+,
\]
where the \(\alpha_i'\) (respectively \(\alpha_j''\)) belong to
\(\Delta'\) (respectively \(\Delta''\)). Applying
\ref{sga3:XXI:3.1.2}, and interchanging \(\Delta'\) and
\(\Delta''\) if necessary, yields a relation
\[
 \delta=\gamma+\beta,
 \qquad
 \gamma\in R_{\Delta'},
 \quad
 \beta\in\Delta'',
 \quad
 \delta\in R.
\]
But \((\beta^\ast,\gamma)=0\), so
\(\gamma-\beta\) is also a root by \ref{sga3:XXI:2.2.5}, which is
impossible.
\end{proof}
